\section{Erd\H{o}s Problem \#50 --- Round 3}

\subsection*{1) ROUND-3 OBJECTIVE}
\textbf{Path (C): obstruction/correction.}
A preliminary correction is necessary: Round~2 carried the wrong problem number. The totient-minimal-preimage question under discussion is Erd\H{o}s Problem~\#50, not \#51; the mathematical content of Round~2 is unchanged, but the numbering is corrected here.
Round~2 left one precise gap: it controlled the size of $n_a/a$ but did \emph{not} control when a candidate $n$ is actually the \emph{minimal} preimage of $a=\varphi(n)$. The most promising Round~3 direction is therefore to attack \emph{minimality} directly.

The new goal is to prove a rigorous obstruction theorem showing that a large class of natural constructions can never yield minimal inverse totients. Concretely, if some subproduct of the prime-power totient blocks of $n$ is one less than a suitable prime, then those blocks can be compressed to a smaller integer with the same totient, so $n$ was not minimal.

\subsection*{2) ROUND-2 FOUNDATION USED}
We use the following previously established items.
\begin{itemize}
\item (R2--Foundation~1, imported from Round~1) If
\[
 n=\prod_{i=1}^r p_i^{\alpha_i},
\]
then
\[
 \varphi(n)=\prod_{i=1}^r p_i^{\alpha_i-1}(p_i-1).
\]
In particular,
\[
 \frac{n}{\varphi(n)}=\prod_{p\mid n}\frac{p}{p-1}.
\]
\item (R2--Foundation~2) The main unresolved issue after Round~2 was the \emph{minimality gap}: even if $n/\varphi(n)$ is large, Round~2 did not provide a mechanism to exclude the possibility that some smaller $m<n$ has $\varphi(m)=\varphi(n)$.
\item (R2--Computation) Up to $n_a\le 10^8$, the record ratio came from
\[
 n_a=3\cdot 5\cdot 17\cdot 47\cdot 347,
\qquad
 a=\varphi(n_a)=2037248,
\qquad
 \frac{n_a}{a}\approx 2.0413788601.
\]
This will be used only as an illustration of the new obstruction.
\end{itemize}

\subsection*{3) NEW INSIGHT / TOOL (ROUND-3)}
The new ingredient is a \textbf{local prime-compression principle}.

Write a candidate preimage in prime-power blocks
\[
 n=\prod_{i=1}^r p_i^{\alpha_i},
 \qquad
 t_i:=\varphi\!\bigl(p_i^{\alpha_i}\bigr)=p_i^{\alpha_i-1}(p_i-1).
\]
Then $\varphi(n)=\prod_i t_i$. If, for some subset of blocks, the product of the corresponding $t_i$'s is of the form $q-1$ with $q$ prime and with $q$ not already forced into the complementary part, then those blocks can be replaced by the single prime $q$.

This produces a smaller integer with the same totient. Hence \emph{minimal} inverse totients must avoid an entire family of such ``Euclid compressions.''

The strongest clean consequence is an \textbf{anti-Euclid condition at the largest prime block}: every subset containing the largest prime-power block gives a product $\prod t_i$ whose successor is composite, except for the tautological singleton corresponding to a largest prime to the first power.

\subsection*{4) ATTACK PLAN (ROUND-3)}
The exact Round~2 gap was: \emph{how can one certify that a high-ratio solution is not minimal?}

The Round~3 plan is:
\begin{enumerate}
\item prove a general compression lemma for prime-power block decompositions;
\item derive from it a clean corollary for every subset containing the largest prime-power block;
\item formulate the resulting obstruction as a necessary condition for any future construction toward Erd\H{o}s~\#50;
\item verify that the Round~2 record-holder satisfies the obstruction.
\end{enumerate}

This overcomes the earlier obstacle because it gives a \emph{direct contradiction to minimality}, not merely a bound on $n/\varphi(n)$.

\subsection*{5) WORK (ROUND-3)}

\subsubsection*{5.1. General compression lemma}
\begin{proposition}[Local prime compression]
Let
\[
 n=\prod_{i=1}^r p_i^{\alpha_i}
\]
with distinct primes $p_i$ and exponents $\alpha_i\ge 1$, and define
\[
 t_i:=\varphi\!\bigl(p_i^{\alpha_i}\bigr)=p_i^{\alpha_i-1}(p_i-1).
\]
Fix a nonempty subset $I\subseteq\{1,\dots,r\}$ and put
\[
 T_I:=\prod_{i\in I} t_i,
 \qquad
 M_I:=\prod_{i\in I} p_i^{\alpha_i}.
\]
Assume that
\begin{enumerate}
\item $q:=T_I+1$ is prime,
\item $q\nmid \prod_{j\notin I} p_j^{\alpha_j}$,
\item either $|I|\ge 2$, or $I=\{i\}$ with $\alpha_i\ge 2$.
\end{enumerate}
Then $n$ is \emph{not} the smallest positive integer $m$ satisfying $\varphi(m)=\varphi(n)$.
\end{proposition}

\begin{proof}
Define
\[
 n':=q\prod_{j\notin I} p_j^{\alpha_j}.
\]
By assumption (2), $q$ is coprime to $\prod_{j\notin I} p_j^{\alpha_j}$, so the multiplicativity of $\varphi$ gives
\[
 \varphi(n')=\varphi(q)\prod_{j\notin I} \varphi\!\bigl(p_j^{\alpha_j}\bigr)
 =(q-1)\prod_{j\notin I} t_j
 =T_I\prod_{j\notin I} t_j
 =\prod_{i=1}^r t_i
 =\varphi(n).
\]
So $n'$ is another preimage of the same totient value.

It remains to show that $n'<n$, equivalently $q<M_I$.

\smallskip
\noindent\emph{Case 1: $I=\{i\}$ and $\alpha_i\ge 2$.}
Then
\[
 q=t_i+1=p_i^{\alpha_i-1}(p_i-1)+1
 =p_i^{\alpha_i}-p_i^{\alpha_i-1}+1
 <p_i^{\alpha_i}=M_I.
\]
Thus $n'<n$.

\smallskip
\noindent\emph{Case 2: $|I|\ge 2$.}
For each $i\in I$,
\[
 t_i=\varphi\!\bigl(p_i^{\alpha_i}\bigr)\le p_i^{\alpha_i}-1,
\]
with equality only when $\alpha_i=1$. Hence
\[
 T_I\le \prod_{i\in I}\bigl(p_i^{\alpha_i}-1\bigr).
\]
Therefore
\[
 q=T_I+1\le \prod_{i\in I}\bigl(p_i^{\alpha_i}-1\bigr)+1.
\]
Since $|I|\ge 2$ and every $p_i^{\alpha_i}\ge 2$, we have the elementary strict inequality
\[
 \prod_{i\in I}(x_i-1)+1<\prod_{i\in I}x_i
 \qquad (x_i\ge 2,\ |I|\ge 2),
\]
obtained by expanding the right-hand side and observing that the omitted positive terms make the difference strictly positive.
Applying this with $x_i=p_i^{\alpha_i}$ yields
\[
 q<\prod_{i\in I} p_i^{\alpha_i}=M_I.
\]
So again $n'<n$.

Thus in all cases $n'$ is a smaller integer with the same totient as $n$, so $n$ cannot be minimal.
\end{proof}

\subsubsection*{5.2. Clean corollary for subsets containing the largest prime block}
The previous proposition has a technical coprimality hypothesis $q\nmid \prod_{j\notin I}p_j^{\alpha_j}$. The next corollary removes that issue by forcing $I$ to contain the largest prime-power block.

\begin{theorem}[Largest-block anti-Euclid obstruction]
Let
\[
 n_a=\prod_{i=1}^r p_i^{\alpha_i}
\qquad (p_1<\cdots <p_r)
\]
be the minimal preimage of a totient value $a$, and put
\[
 t_i:=p_i^{\alpha_i-1}(p_i-1).
\]
Let $I\subseteq\{1,\dots,r\}$ be any subset with $r\in I$.
If $I\neq \{r\}$ or $\alpha_r\ge 2$, then
\[
 1+\prod_{i\in I} t_i
\]
is composite.
Equivalently: the only subset containing the largest prime-power block whose $t$-product can have prime successor is the tautological singleton $I=\{r\}$ when $\alpha_r=1$, where the successor is just $p_r$ itself.
\end{theorem}

\begin{proof}
Suppose, toward a contradiction, that for some such subset $I$ the number
\[
 q:=1+\prod_{i\in I}t_i
\]
is prime.
We verify the hypotheses of the local compression proposition.

First, $q$ does not divide the complementary factor $\prod_{j\notin I}p_j^{\alpha_j}$. Indeed, all prime divisors in the complement are among $p_1,\dots,p_{r-1}$, hence are $<p_r$. On the other hand,
\[
 q=1+\prod_{i\in I} t_i \ge 1+t_r.
\]
If $\alpha_r\ge 2$, then
\[
 1+t_r=1+p_r^{\alpha_r-1}(p_r-1)>p_r.
\]
If $\alpha_r=1$, then $t_r=p_r-1$, so $q\ge p_r$, with equality only if every additional factor in the product equals $1$; this still causes no problem because $p_r\notin\{p_j:j\notin I\}$ when $r\in I$.
Thus in every case $q$ is not a prime factor of the complement.

Second, the hypothesis ``either $|I|\ge 2$ or the relevant exponent is at least $2$'' is exactly the present assumption that $I\neq \{r\}$ or $\alpha_r\ge 2$.

Therefore Proposition~5.1 applies and shows that $n_a$ is not minimal, a contradiction.
Hence $q$ must be composite.
\end{proof}

\subsubsection*{5.3. Immediate consequences}
\begin{corollary}[Largest prime power cannot compress]
If $p_r^{\alpha_r}\parallel n_a$ with $\alpha_r\ge 2$, then
\[
 \varphi\!\bigl(p_r^{\alpha_r}\bigr)+1
 =p_r^{\alpha_r-1}(p_r-1)+1
\]
is composite.
\end{corollary}

\begin{proof}
Take $I=\{r\}$ in Theorem~5.2.
\end{proof}

\begin{corollary}[Largest prime with one additional block]
Assume $n_a$ is squarefree and write
\[
 n_a=p_1\cdots p_r
\qquad (p_1<\cdots <p_r).
\]
Then for every $j<r$,
\[
 (p_j-1)(p_r-1)+1
\]
is composite.
More generally, for every nonempty subset $J\subseteq\{1,\dots,r-1\}$,
\[
 1+(p_r-1)\prod_{j\in J}(p_j-1)
\]
is composite.
\end{corollary}

\begin{proof}
Apply Theorem~5.2 with $I=J\cup\{r\}$ and note that $t_i=p_i-1$ in the squarefree case.
\end{proof}

\begin{corollary}[Fixed-base obstruction]
Fix an odd integer
\[
 m=\prod_{j=1}^s r_j^{\beta_j}
\]
with distinct prime bases $r_j$, and let
\[
 u_J:=\prod_{j\in J}\varphi\!\bigl(r_j^{\beta_j}\bigr)
\qquad (\emptyset\neq J\subseteq\{1,\dots,s\}).
\]
Let $q$ be a prime exceeding every prime divisor of $m$.
If $n=mq$ is the minimal preimage of $\varphi(n)$, then for every nonempty $J$,
\[
 u_J(q-1)+1
\]
is composite.
\end{corollary}

\begin{proof}
Write the prime-power blocks of $m$ together with the final block $q$.
Since the $q$-block is the largest prime block and appears to the first power, apply Theorem~5.2 to subsets consisting of that block together with the blocks indexed by $J$.
The corresponding product of $t$-values is exactly $u_J(q-1)$.
\end{proof}

\subsubsection*{5.4. How this sharpens the Round-2 record-holder}
Round~2 found the record-holder up to $10^8$:
\[
 n_a=3\cdot 5\cdot 17\cdot 47\cdot 347,
\qquad
 a=\varphi(n_a)=2037248.
\]
Here the largest prime block is $347$, and the smaller $t$-values are
\[
 2,\ 4,\ 16,\ 46.
\]
Thus Theorem~5.2 forces the compositeness of
\[
 1+346u
\]
for every nonempty subset-product
\[
 u\in\{2,4,8,16,32,46,64,92,128,184,368,736,1472,2944,5888\}.
\]
A direct factorization check gives, for example,
\[
 693=3^2\cdot 7\cdot 11,
\qquad
 15917=11\cdot 1447,
\qquad
 2037249=3^2\cdot 41\cdot 5521,
\]
and similarly for the remaining $12$ values.

So the Round~2 record-holder already exhibits the new Round~3 obstruction in a strong form: all $15$ nontrivial largest-block Euclid successors are composite.

\subsubsection*{5.5. The corrected obstruction problem}
Theorem~5.2 isolates a precise barrier that any future proof of Erd\H{o}s~\#50 must overcome.

\begin{quote}
\emph{Corrected obstruction statement.}
Any infinite family of minimal inverse totients
\[
 n_\nu=\prod_{i=1}^{r_\nu} p_{\nu,i}^{\alpha_{\nu,i}}
\]
can witness $n_{a_\nu}/a_\nu\to\infty$ only if, for every $\nu$ and every subset containing the largest prime-power block, the successor
\[
 1+\prod_{i\in I} \varphi\!\bigl(p_{\nu,i}^{\alpha_{\nu,i}}\bigr)
\]
is composite, except for the tautological singleton largest-prime block.
\end{quote}

This is a genuine obstruction: it rules out any construction strategy in which these Euclid-type successors are allowed to be prime infinitely often.

\subsection*{6) ADVERSARIAL VERIFICATION}
\begin{itemize}
\item \textbf{Why the unrestricted statement is false.}
A too-optimistic claim would say that for \emph{every} nontrivial subset $I$, the successor $1+\prod_{i\in I}t_i$ must be composite. This is false because the new prime can coincide with a prime already present in the complementary part.
For example,
\[
 n=3\cdot 7\cdot 13
\]
has
\[
 (3-1)(7-1)+1=13,
\]
which is prime, yet this does not by itself yield a coprime replacement because $13$ already occurs outside the chosen subset. This is exactly why Proposition~5.1 includes the coprimality hypothesis, and why Theorem~5.2 is stated only for subsets containing the largest prime block.

\item \textbf{Boundary case $I=\{r\}$ with $\alpha_r=1$.}
Then
\[
 1+t_r=1+(p_r-1)=p_r
\]
is prime, but replacing the largest prime block by itself changes nothing. This tautological case is explicitly excluded from Theorem~5.2.

\item \textbf{Strict inequality $q<M_I$.}
This was checked separately in the proof of Proposition~5.1.
For singleton subsets this requires $\alpha_i\ge 2$; otherwise $q=p_i$ and no decrease occurs.
For subsets of size at least $2$, the inequality
\[
 \prod (x_i-1)+1<\prod x_i
\]
for $x_i\ge 2$ is strict and leaves no hidden equality case.

\item \textbf{Interaction with Round~2 computation.}
The theorem does not assert that the Round~2 record-holder is minimal; that was already known computationally from Round~2. What Round~3 adds is a structural explanation of why certain obvious compressions are impossible: all $15$ relevant largest-block Euclid successors are indeed composite.
\end{itemize}

\subsection*{7) FINAL}
\textbf{UNRESOLVED (BUT STRICTLY ADVANCED).}

Round~3 resolves the main logical weakness left after Round~2 by proving a new \emph{minimality obstruction}:
\begin{itemize}
\item the general local prime-compression lemma (Proposition~5.1), and
\item the largest-block anti-Euclid obstruction (Theorem~5.2).
\end{itemize}
These results rule out an entire class of natural construction strategies. In particular, any attempt to prove Erd\H{o}s~\#50 by exhibiting explicit high-ratio families must now solve a sharply formulated compositeness problem for all Euclid-type successors attached to the largest prime-power block.

This is strictly stronger than Round~2: it does not merely bound or search for $n_a/a$, but gives a concrete criterion showing when a candidate preimage cannot be minimal.

\subsection*{8) COMPLETION ESTIMATE (MANDATORY)}
\textbf{COMPLETION: 48\%}

\subsection*{9) REFERENCES}
\begin{itemize}
\item T. F. Bloom, \emph{Erd\H{o}s Problem \#50}, \texttt{https://www.erdosproblems.com/50}, accessed 2026-03-07.
\item J. B. Rosser and L. Schoenfeld, ``Approximate formulas for some functions of prime numbers'', \emph{Illinois J. Math.} \textbf{6} (1962), Theorem~15.
\item E. W. Weisstein, ``Totient Function'', MathWorld.
\end{itemize}
